\section{Core Idea}

Let $dp[i][j]$ be the minimum number of edits needed to transform the first
$i$ characters of $a$ into the first $j$ characters of $b$.

The last operation must be exactly one of these three:

\begin{itemize}
\item delete $a[i-1]$
\item insert $b[j-1]$
\item align $a[i-1]$ with $b[j-1]$, which costs $0$ if they match and $1$ if we replace
\end{itemize}

So we get

$$
dp[i][j] =
\min\left(
dp[i-1][j] + 1,\;
dp[i][j-1] + 1,\;
dp[i-1][j-1] + c(i,j)
\right),
$$

where $c(i,j)=0$ if $a[i-1]=b[j-1]$ and $1$ otherwise.

The base cases are $dp[i][0] = i$ and $dp[0][j] = j$.

\section{Toy Example}

For the animation, we use

\begin{itemize}
\item $a = \texttt{CAT}$
\item $b = \texttt{CUTE}$
\end{itemize}

This example is small enough to fit on screen, but it still shows the three
important ideas:

\begin{itemize}
\item a match with zero extra cost
\item a replacement with cost $1$
\item an insertion that increases the answer at the end
\end{itemize}

\section{Walkthrough}

\begin{animation}[id="edit-distance-dp", label="Edit Distance DP walkthrough"]
\shape{dp}{DPTable}{rows=4, cols=5, label="dp[i][j]"}
\shape{vars}{VariableWatch}{names=["a","b","cell","rule"], label="Current state"}
\apply{vars.var[a]}{value="CAT"}
\apply{vars.var[b]}{value="CUTE"}
\apply{vars.var[cell]}{value="start"}
\apply{vars.var[rule]}{value="set up states"}

\step
\recolor{dp.cell[0][0]}{state=current}
\narrate{Row $i$ means the first $i$ characters of \texttt{CAT}, and column $j$ means the first $j$ characters of \texttt{CUTE}. Each cell stores the best answer for those two prefixes.}

\step
\apply{dp.cell[0][0]}{value="0"}
\apply{dp.cell[0][1]}{value="1"}
\apply{dp.cell[0][2]}{value="2"}
\apply{dp.cell[0][3]}{value="3"}
\apply{dp.cell[0][4]}{value="4"}
\apply{dp.cell[1][0]}{value="1"}
\apply{dp.cell[2][0]}{value="2"}
\apply{dp.cell[3][0]}{value="3"}
\recolor{dp.cell[0][0]}{state=done}
\recolor{dp.cell[0][1]}{state=done}
\recolor{dp.cell[0][2]}{state=done}
\recolor{dp.cell[0][3]}{state=done}
\recolor{dp.cell[0][4]}{state=done}
\recolor{dp.cell[1][0]}{state=done}
\recolor{dp.cell[2][0]}{state=done}
\recolor{dp.cell[3][0]}{state=done}
\apply{vars.var[cell]}{value="base row and column"}
\apply{vars.var[rule]}{value="only insertions or deletions"}
\narrate{The first row and first column are immediate. Reaching the empty string uses only deletions, and building a prefix from the empty string uses only insertions.}

\step
\apply{dp.cell[1][1]}{value="0"}
\apply{dp.cell[1][2]}{value="1"}
\apply{dp.cell[2][1]}{value="1"}
\recolor{dp.cell[1][1]}{state=current}
\recolor{dp.cell[1][2]}{state=done}
\recolor{dp.cell[2][1]}{state=done}
\annotate{dp.cell[1][1]}{label="match C", arrow_from="dp.cell[0][0]", color=good}
\apply{vars.var[cell]}{value="dp[1][1]"}
\apply{vars.var[rule]}{value="matching characters use the diagonal"}
\narrate{The first real transition is $dp[1][1]$. Since \texttt{C} matches \texttt{C}, the diagonal move costs $0$, so $dp[1][1] = dp[0][0] = 0$.}

\step
\apply{dp.cell[2][2]}{value="1"}
\recolor{dp.cell[1][1]}{state=good}
\recolor{dp.cell[1][2]}{state=done}
\recolor{dp.cell[2][1]}{state=done}
\recolor{dp.cell[2][2]}{state=current}
\annotate{dp.cell[2][2]}{label="delete", arrow_from="dp.cell[1][2]", color=info}
\annotate{dp.cell[2][2]}{label="insert", arrow_from="dp.cell[2][1]", color=info}
\annotate{dp.cell[2][2]}{label="replace", arrow_from="dp.cell[1][1]", color=good}
\apply{vars.var[cell]}{value="dp[2][2]"}
\apply{vars.var[rule]}{value="take the minimum of three choices"}
\narrate{Now compare \texttt{CA} with \texttt{CU}. Deleting and inserting both give cost $2$, while replacing \texttt{A} with \texttt{U} gives $dp[1][1] + 1 = 1$. So the minimum is $1$.}

\step
\apply{dp.cell[2][3]}{value="2"}
\apply{dp.cell[3][3]}{value="1"}
\apply{dp.cell[3][4]}{value="2"}
\recolor{dp.cell[2][2]}{state=path}
\recolor{dp.cell[3][3]}{state=path}
\recolor{dp.cell[3][4]}{state=good}
\annotate{dp.cell[3][3]}{label="match T", arrow_from="dp.cell[2][2]", color=good}
\annotate{dp.cell[3][4]}{label="insert E", arrow_from="dp.cell[3][3]", color=path}
\apply{vars.var[cell]}{value="dp[3][4]"}
\apply{vars.var[rule]}{value="finish with one insertion"}
\narrate{The rest of the table is filled the same way. Here \texttt{T} matches \texttt{T}, so $dp[3][3] = 1$, and then we insert the final \texttt{E}, giving $dp[3][4] = 2$. The answer is the bottom-right cell.}
\end{animation}

\section{Proof Sketch}

Consider an optimal edit sequence for the prefixes $a[0..i-1]$ and $b[0..j-1]$.
Its final step must be one of three cases:

\begin{itemize}
\item If it deletes $a[i-1]$, then before that step we must already have
transformed $a[0..i-2]$ into $b[0..j-1]$, so the cost is $dp[i-1][j] + 1$.
\item If it inserts $b[j-1]$, then before that step we must already have
transformed $a[0..i-1]$ into $b[0..j-2]$, so the cost is $dp[i][j-1] + 1$.
\item Otherwise it aligns the last characters of both prefixes. If they match,
the extra cost is $0$; otherwise we pay $1$ for a replacement. So this case
costs $dp[i-1][j-1] + c(i,j)$.
\end{itemize}

Every optimal solution ends in one of these cases, and each case reduces the
problem to a smaller prefix pair. Therefore taking the minimum of the three
values gives the correct answer for $dp[i][j]$.

\section{Complexity}

There are $(n+1)(m+1)$ states, and each state is computed in constant time.

\begin{itemize}
\item Time complexity: $O(nm)$
\item Memory complexity: $O(nm)$
\end{itemize}

\section{Pitfalls}

\begin{itemize}
\item The table size is $(n+1) \times (m+1)$ because we need the empty-prefix
row and column.
\item The transition uses $a[i-1]$ and $b[j-1]$ because $i$ and $j$ are prefix
lengths, not direct indices into the DP table.
\item For the original constraints, a full 2D table is acceptable with
\texttt{int}, but unnecessary larger types waste memory.
\end{itemize}

\section{Reference Implementation}

\begin{lstlisting}[language=C++]
#include <algorithm>
#include <iostream>
#include <vector>
using namespace std;

int main() {
    string a, b;
    cin >> a >> b;

    int n = a.size();
    int m = b.size();

    vector<vector<int>> dp(n + 1, vector<int>(m + 1));

    for (int i = 0; i <= n; ++i) {
        dp[i][0] = i;
    }
    for (int j = 0; j <= m; ++j) {
        dp[0][j] = j;
    }

    for (int i = 1; i <= n; ++i) {
        for (int j = 1; j <= m; ++j) {
            int cost = (a[i - 1] == b[j - 1] ? 0 : 1);
            dp[i][j] = min({
                dp[i - 1][j] + 1,
                dp[i][j - 1] + 1,
                dp[i - 1][j - 1] + cost,
            });
        }
    }

    cout << dp[n][m] << "\n";
}
\end{lstlisting}
